\documentclass[DIV=calc]{scrartcl}
\usepackage[utf8]{inputenc}
\usepackage[T1]{fontenc}
\usepackage[ngerman]{babel}
\usepackage{graphicx}

\usepackage{fixltx2e}
\usepackage{ellipsis}
\usepackage[tracking=true]{microtype}

\usepackage{lmodern}      % Ersatz fuer Computer Modern-Schriften
\usepackage{hfoldsty}

\usepackage{fourier} 			% Schriftart
\usepackage[scaled=0.81]{helvet} 	% Schriftart

\usepackage{url}
\usepackage{tocloft} 			% Paket für Table of Contents

\usepackage{xcolor}
\definecolor{urlred}{HTML}{660000}

\usepackage{hyperref}
\hypersetup{
 colorlinks=true,	
 linkcolor=black,	% Farbe der internen Links (u.a. Table of Contents)
 urlcolor=black,	% Farbe der url-links
 citecolor=black} % Farbe der Literaturverzeichnis-Links

\usepackage{mdwlist} 	% Änderung der Zeilenabstände bei itemize und enumerate

\parindent 0pt 				% Absatzeinrücken verhindern
\parskip 12pt 				% Absätze durch Lücke trennen

%\usepackage{titlesec}	% Abstand nach Überschriften neu definieren
%\titlespacing{\subsection}{0ex}{3ex}{-1ex}
%\titlespacing{\subsubsection}{0ex}{3ex}{-1ex}		

% \pagestyle{empty}
\setlength{\textheight}{23cm}
\usepackage{fancyhdr}
\pagestyle{fancy}
\cfoot{}
\lfoot{Zusammenkunft aller Physik-Fachschaften}
\rfoot{www.zapfev.de\\stapf@zapf.in}
\renewcommand{\headrulewidth}{0pt}
\renewcommand{\footrulewidth}{0.1pt}


\begin{document}
\hspace{0.87\textwidth}
\begin{minipage}{120pt}
\vspace{-1.8cm}
\includegraphics[width=80pt]{logo.pdf}
\centering
\small Zusammenkunft aller Physik-Fachschaften
\end{minipage}
\begin{center}
\huge{Resolution der Zusammenkunft aller Physik-Fachschaften} \\
\normalsize
\end{center}

\vspace{1cm}
\section*{Resolution zu den geplanten Studiengebühren in Baden-Württemberg}

Die Landesregierung Baden-Württemberg hat vor, Studiengebühren für Nicht-EU-Ausländer in Höhe von 1.500 € pro Semester, für Zweitstudierende in Höhe von 650 € pro Semester einzuführen, sowie die Verwaltungskosten von 60 € auf 70 € zu erhöhen um die Lücken im Landeshaushalt zu schließen. \\
Wir, die Zusammenkunft aller deutschsprachigen Physikfachschaften, unterstützen die Landesstudierendenvertretung Baden-Württemberg in ihrer Positionierung\footnote{\url{http://www.studis.de/lak-bawue/fileadmin/lak-bawue/PMs_und_offene_Briefe/Positionierung_der_Landesstudierendenvertretung_Baden-Wuerttemberg_zur_Einfuehrung_von_Studiengebuehren.pdf}} gegen dieses Vorhaben und sprechen uns gegen Studiengebühren jeglicher Form aus.\\
Aus gegebenem Anlass möchten wir aus dem Parteitagsbeschluss\footnote{\href{http://www.gruene.de/fileadmin/user_upload/Beschluesse/Bildung-keine-Studiengebuehren-Beschluss-BDK-Erfurt-11-2008.pdf}{\texttt{http://www.gruene.de/fileadmin/user\_upload/Beschluesse/Bildung-keine- Studiengebuehren-Beschluss-BDK-Erfurt-11-2008.pdf}}} der GRÜNEN vom November 2008 mit dem Titel: „Grüne für klares Nein zu Studiengebühren“ zitieren: „Studiengebühren führen in eine hochschulpolitische Sackgasse“. „Die Unterfinanzierung […] muss über andere Wege überwunden werden, Studiengebühren sind das falsche Mittel dafür.“\\





\vfill
\begin{flushright}
Verabschiedet am 12.11.2016 in Dresden
\end{flushright}




\end{document}
